% !PRINTED 207
\weberpage{207}

%<<CONT
sets. Formula (2) of § 21 then gives:

\weq{6}{C = (-1)^{(g_1+1)(g_2+1)\frac{n-1}{2}} n \frac{\vartheta_{g_1,g_2}(0,\omega)}{\vartheta_{g_1,g_2}(2\nu'\omega, n\omega)} \frac{\vartheta_{11}'(0, n\omega)}{\vartheta_{11}'(0, \omega)}.}

One obtains the \gesperrt{Abelian} relations quite simply by setting
$nu = 2\nu'\omega$ in (5), whereby the right-hand side vanishes:

\weq{7}{\sum^{\nu} e^{\frac{2\nu\nu'\pi i}{n}} \frac{\vartheta_{11}\left(\frac{2\nu'\omega + 2\nu}{n}\right)}{\vartheta_{g_1,g_2}\left(\frac{2\nu'\omega + 2\nu}{n}\right)} = 0.}
% [sic] The German source prints 2 in this exponential, although direct
% substitution in (5), and the downstream formulae (8)--(11), retain 8.

If one applies to the $\vartheta$-functions occurring here a
linear transformation

\[
\begin{pmatrix} \alpha, & \beta \\ \gamma, & \delta \end{pmatrix}, \qquad \alpha\delta - \beta\gamma = 1
\]

then by § 39 one obtains the general formula:

\weq{8}{\sum^{\nu} e^{\frac{8\nu\nu'\pi i}{n}} \frac{\vartheta_{11}\left(\frac{2(\nu\alpha+\nu'\gamma) + 2(\nu\beta+\nu'\delta)\omega}{n}\right)}{\vartheta_{g_1,g_2}\left(\frac{2(\nu\alpha+\nu'\gamma) + 2(\nu\beta+\nu'\delta)\omega}{n}\right)} = 0.}

These equations now pass over, when one sets $g_1, g_2 = 0, 1$
and introduces the notation of § 42, directly into relations between the roots $x_{\nu,\nu'}$ (§ 61):

\weq{9}{\sum^{\nu} e^{\frac{8\nu\nu'\pi i}{n}} x_{\nu\alpha+\nu'\gamma,\ \nu\beta+\nu'\delta} = 0,}

and to the two special cases:

\[
\begin{pmatrix} \alpha, & \beta \\ \gamma, & \delta \end{pmatrix} = \begin{pmatrix} 1, & 0 \\ 0, & 1 \end{pmatrix}, \qquad \begin{pmatrix} 0, & 1 \\ -1, & 0 \end{pmatrix}
\]

correspondingly:

\weq{10}{\sum^{\nu} e^{\frac{8\nu\nu'\pi i}{n}} x_{\nu,\nu'} = 0,}

\weq{11}{\sum^{\nu'} e^{-\frac{8\nu\nu'\pi i}{n}} x_{\nu,\nu'} = 0.}

By (5), (6) one can also determine the value of these sums, if in them the $n$th root of unity

\[
e^{\frac{8\pi i}{n}}
\]
%CONT>>
